\section*{Capítulo \ref{ch:g8:plastics} --- Os plásticos e os materiais sintéticos}

\begin{solution}{exo:g8:plastics:1}
Naturais: algodão, lã. Artificiais: papel, viscose. Sintéticos: náilon,
poliestireno.
\end{solution}

\begin{solution}{exo:g8:plastics:2}
Do petróleo (ou do gás natural).
\end{solution}

\begin{solution}{exo:g8:plastics:3}
O PET (poli(tereftalato de etileno)): garrafas de água e de refrigerante.
\end{solution}

\begin{solution}{exo:g8:plastics:4}
Um \omterm{def:g8:plastics:thermoplastic}{termoplástico} amolece cada vez que é aquecido e pode ser moldado de
novo; um \omterm{def:g8:plastics:thermoplastic}{termofixo} endurece para sempre da primeira vez e carboniza em
vez de amolecer.
\end{solution}

\begin{solution}{exo:g8:plastics:5}
O polipropileno (PP), um \omterm{def:g8:plastics:thermoplastic}{termoplástico}: ele pode ser derretido e moldado de novo.
\end{solution}

\begin{solution}{exo:g8:plastics:6}
O triângulo só dá o nome do \omterm{def:g8:plastics:plastic}{plástico}, para ajudar na separação. Se o
objeto vai ser reciclado depende de ele ser recolhido e de haver uma
usina que recicle esse \omterm{def:g8:plastics:plastic}{plástico}.
\end{solution}

\begin{solution}{exo:g8:plastics:7}
Uma frigideira esquenta: um \omterm{def:g8:plastics:thermoplastic}{termoplástico} amoleceria, um \omterm{def:g8:plastics:thermoplastic}{termofixo} não.
\end{solution}

\begin{solution}{exo:g8:plastics:8}
$\frac{353}{460} \approx 0.77$: cerca de \qty{77}{\percent}.
\end{solution}

\begin{solution}{exo:g8:plastics:9}
O PVC contém cloro: queimá-lo libera cloreto de hidrogênio, um gás
corrosivo, além da fuligem e do monóxido de carbono da \omterm{def:g8:combustion-and-fuels:complete}{combustão
incompleta} de uma fogueira de quintal.
\end{solution}

\begin{solution}{exo:g8:plastics:10}
$300 \times 36 = 10\,800$ garrafas; $10\,800 \times 25 = 270\,000$\,g,
isto é, \qty{270}{kg}.
\end{solution}

\begin{solution}{exo:g8:plastics:11}
Dióxido de carbono e água. O carbono dele vem do petróleo, guardado no
subsolo durante milhões de anos: queimá-lo acrescenta dióxido de carbono
ao \omterm{def:g4:air-a-mixture-of-gases:air}{ar}, como faz a \omterm{def:g4:what-a-fire-needs:burning}{queima} de um \omterm{def:g8:combustion-and-fuels:fossil}{combustível fóssil}.
\end{solution}

\begin{solution}{exo:g8:plastics:12}
A garrafa não apodrece: o sol e as ondas a quebram em pedaços cada vez
menores. Os pedacinhos boiam ou afundam na água, e os peixes os engolem
junto com o alimento.
\end{solution}

\begin{solution}{pb:g8:plastics:1}
\textbf{1.} Corpos: 1 (PET). Tampas: 2 (PEAD) ou 5 (PP).

\textbf{2.} Sintéticos: eles são produzidos a partir de substâncias
construídas pelos químicos, principalmente a partir do petróleo.

\textbf{3.} Sim: o PET é um \omterm{def:g8:plastics:thermoplastic}{termoplástico}.

\textbf{4.} Cada \omterm{def:g8:plastics:plastic}{plástico} é reciclado separadamente: derretida junta,
uma \omterm{def:g6:pure-substances-and-mixtures:mixture}{mistura} de \omterm{def:g8:plastics:plastic}{plásticos} dá um material de má qualidade.

\textbf{5.} $0.80 \times 1000 = \qty{800}{kg}$.

\textbf{6.} Tampas: $0.12 \times 1000 = \qty{120}{kg}$. Rótulos: $0.05
\times 1000 = \qty{50}{kg}$.

\textbf{7.} $0.03 \times 1000 = \qty{30}{kg}$ de sujeira e líquido.

\textbf{8.} $80 + 12 + 5 = \qty{97}{\percent}$.

\textbf{9.} $0.025 \times 800 = \qty{20}{kg}$.

\textbf{10.} Os flocos poupam petróleo e a energia de produzir PET novo,
e transformam lixo num objeto útil.

\textbf{11.} O dióxido de carbono (e, se a \omterm{def:g4:what-a-fire-needs:burning}{queima} for ruim, o monóxido
de carbono e a fuligem).

\textbf{12.} $800 - 20 = \qty{780}{kg}$ de flocos de PET limpos por fardo.
\end{solution}
